% answers_short.tex —— 2013 年·安徽卷·理科 · 简短答案速查
% 由 tools/make_short_answers.py 自 answers.tex 的 \gksolhead/\gkanswer 提取，
% 并与 2013/2013_数学真题_安徽理_答案要点.md 逐题比对后生成，勿手改（改源 MD 后重跑）。

\begin{gknotice}
\textbf{使用说明}：本节只给最终答案便于快速核对，完整推导见解析版。答案为编写者独立推导，\textbf{非官方标准答案}。
\end{gknotice}

\gkpart{选择题}
\begin{center}\begin{tabular}{ccccccccccc}
\toprule
题号 & 1 & 2 & 3 & 4 & 5 & 6 & 7 & 8 & 9 & 10 \\
\midrule
答案 & A & D & A & C & C & D & B & B & D & A \\
\bottomrule
\end{tabular}\end{center}

\gkpart{填空题}
\begin{itemize}[leftmargin=2.6em]
  \item[\textbf{11.}] $\dfrac12$。$x^4$ 项对应 $k=3$，由 $56a^3=7$ 得 $a=\dfrac12$。
  \item[\textbf{12.}] $\dfrac{2\pi}{3}$。由 $3a=5b$ 与 $b+c=2a$ 得 $b=\dfrac35a$，$c=\dfrac75a$，故 $\cos C=-\dfrac12$。
  \item[\textbf{13.}] $[1,+\infty)$。设 $C(t,t^2)$，由 $\overrightarrow{CA}\cdot\overrightarrow{CB}=0$ 得 $(t^2-a)(t^2-a+1)=0$，非退化情形须 $a-1\geqslant0$。
  \item[\textbf{14.}] $\sqrt{3n-2}$。梯形面积与 $a_{n+1}^2-a_n^2$ 成正比且为常数，故 $\{a_n^2\}$ 是公差为 $3$ 的等差数列。
  \item[\textbf{15.}] ①②③⑤。建系得截面方程 $tx-2ty+z=0$，按 $t=CQ$ 分类：$0<t<\dfrac12$ 为四边形，$t=\dfrac12$ 为等腰梯形，$t=\dfrac34$ 时 $C_1R=\dfrac13$，$\dfrac12<t<1$ 为五边形，$t=1$ 时面积为 $\dfrac{\sqrt6}{2}$。
\end{itemize}

\gkpart{解答题}
\begin{itemize}[leftmargin=2.6em]
  \item[\textbf{16.}] (1) $\omega=1$。$f(x)=2\sin\left(2\omega x+\dfrac\pi4\right)+\sqrt2$，由 $\dfrac{\pi}{\omega}=\pi$ 得 $\omega=1$；(2) 在 $\left[0,\dfrac{\pi}{8}\right]$ 上单调递增，在 $\left[\dfrac{\pi}{8},\dfrac{\pi}{2}\right]$ 上单调递减。由 $2x+\dfrac\pi4\in\left[\dfrac\pi4,\dfrac{5\pi}4\right]$ 判 $\sin$ 的单调性即得
  \item[\textbf{17.}] (1) $\dfrac{a}{1+a^2}$。$I=\left(0,\dfrac{a}{1+a^2}\right)$，长度即右端点；(2) $\dfrac{1-k}{k^2-2k+2}$。$L(a)=\dfrac{a}{1+a^2}$ 在 $(0,1]$ 增、$[1,+\infty)$ 减，且 $L(1-k)<L(1+k)$，故最小值在 $a=1-k$ 处取得
  \item[\textbf{18.}] (1) $\dfrac{8x^2}{5}+\dfrac{8y^2}{3}=1$。$c^2=2a^2-1$，由 $2c=1$ 得 $a^2=\dfrac58$，$b^2=\dfrac38$；(2) $x+y=1$。由 $F_1P\perp F_1Q$ 得 $x_0^2-y_0^2=c^2$，联立椭圆知 $P(a^2,1-a^2)$，故 $x_P+y_P=1$
  \item[\textbf{19.}] (1) 见解析版；(2) $17-12\sqrt2$。设高 $h$、底面半径 $R$，$h=R\tan22.5^\circ$；由 $OP$ 与平面 $PCD$ 成 $60^\circ$ 得弦心距 $d=\sqrt3h$，故 $\cos\angle COD=2\left(\dfrac dR\right)^2-1=17-12\sqrt2$
  \item[\textbf{20.}] (1) 见解析版；(2) 见解析版
  \item[\textbf{21.}] (1) $\dfrac{k(2n-k)}{n^2}$。甲未收到任一位老师通知的概率为 $\left(1-\dfrac kn\right)^2$，用对立事件即得；(2) 令 $u_0=\dfrac{k(n-k)-1}{n+2}$。当 $u_0\notin\mathbf Z$ 时，$m=k+\lceil u_0\rceil$；当 $u_0\in\mathbf Z$ 时，$m=k+u_0$ 或 $m=k+u_0+1$。由 $\dfrac{a_{u+1}}{a_u}=\dfrac{(k-u)(n-k-u)}{(u+1)^2}$ 与 $1$ 比较即得
\end{itemize}
